% ROZDZIAŁ 25 — RÓŻNICZKOWANIE CAŁEK I MIAR
\chapter{Różniczkowanie całek i miar}
\label{ch:differentiation-25}

\section{Dlaczego różniczkowanie jest twierdzeniem o średnich?}
Niech $d\ge1$, a $\lambda_d$ oznacza miarę Lebesgue'a na
$\mathbb R^d$. Dla $x\in\mathbb R^d$ i $r>0$ piszemy
$B(x,r)=\{y:|y-x|<r\}$ oraz $v_d=\lambda_d(B(0,1))$.
Jeżeli $f\in L^1_{\mathrm{loc}}(\mathbb R^d)$, to na pierwszy
rzut oka naturalne jest oczekiwanie, że średnia $f$ po kuli
malejącej do punktu $x$ dąży do $f(x)$. Nie zachodzi to
we wszystkich punktach ani dla dowolnego przedstawiciela
funkcji. Poniższa teoria precyzyjnie określa, w jakim
sensie odtwarzamy funkcję z jej całek.

\begin{definicja}[Wartość średnia i punkt Lebesgue'a]
\label{def:lebesgue-point-25}
Dla $f\in L^1_{\mathrm{loc}}(\mathbb R^d)$ definiujemy
$A_rf(x)=\lambda_d(B(x,r))^{-1}\int_{B(x,r)}f(y)\,\mathrm d\lambda_d(y)$.
Punkt $x$ nazywamy \emph{punktem Lebesgue'a} wybranego
przedstawiciela $f$, jeżeli
\begin{equation}\label{eq:lebesgue-point-25}
 \lim_{r\downarrow0}\frac1{v_dr^d}
 \int_{B(x,r)}|f(y)-f(x)|\,\mathrm d\lambda_d(y)=0.
\end{equation}
W szczególności $A_rf(x)\to f(x)$ w takim punkcie.
Przypominamy, że wartości $f(x)$ na zbiorze miary zero
można zmieniać bez zmiany klasy elementu $L^1_{\mathrm{loc}}$.
\end{definicja}

\begin{przyklad}[Nieciągłość w zbiorze zerowym]
Jeżeli $f=\chi_{\mathbb Q^d}$, to $A_rf(x)=0$ dla
wszystkich $x,r$, ponieważ $\lambda_d(\mathbb Q^d)=0$.
Żaden punkt wymierny nie jest punktem Lebesgue'a tego
konkretnego przedstawiciela, a każdy niewymierny jest.
Równość prawie wszędzie nie może więc zostać zastąpiona
równością w każdym punkcie.
\end{przyklad}

\section{Wzmocnienie nierówności maksymalnej}
\begin{definicja}[Maksymalna funkcja miary]
\label{def:max-measure-25}
Jeżeli $\nu$ jest skończoną dodatnią miarą borelowską
na $\mathbb R^d$, określamy
$M\nu(x)=\sup_{r>0}\nu(B(x,r))/(v_dr^d)$.
Dla miary ze znakiem bądź zespolonej podstawiamy jej
wariację całkowitą $|\nu|$.
\end{definicja}

\begin{twierdzenie}[Oszacowanie słabego typu dla miar]
\label{thm:weak-measure-25}
Niech $\nu$ będzie skończoną dodatnią miarą borelowską
na $\mathbb R^d$. Dla każdego $t>0$ zachodzi
\begin{equation}\label{eq:weak-measure-25}
 \lambda_d\bigl(\{x:M\nu(x)>t\}\bigr)
 \le \frac{5^d}{t}\nu(\mathbb R^d).
\end{equation}
To samo oszacowanie obowiązuje z $\nu$ zastąpioną przez
$|\nu|$, gdy $\nu$ jest miarą ze znakiem lub zespoloną.
\end{twierdzenie}
\begin{proof}
Dla ustalonego $r>0$ funkcja $x\mapsto\nu(B(x,r))$
jest półciągła z dołu: jeśli $x_n\to x$, to dla każdego
$y\in B(x,r)$ ostatecznie $y\in B(x_n,r)$, zatem lemat
Fatou zastosowany do miary $\nu$ daje
$\nu(B(x,r))\le\liminf_n\nu(B(x_n,r))$. Stąd
$\{M\nu>t\}$ jest otwarty. Niech $K$ będzie jego
zwartym podzbiorem. Dla każdego $x\in K$ istnieje
kula $B_x=B(x,r_x)$ taka, że
$\nu(B_x)>tv_dr_x^d$. Ponieważ $\nu$ jest skończona,
promienie są ograniczone z góry. Z lematu $5r$
(twierdzenie~\ref{thm:5r-24}) wybieramy rozłączne
kule $B_j$ o pokryciu $K\subseteq\bigcup_j5B_j$.
Wtedy
\[
\lambda_d(K)\le5^d\sum_jv_dr_j^d
 <\frac{5^d}{t}\sum_j\nu(B_j)
 \le\frac{5^d}{t}\nu(\mathbb R^d).
\]
Korzystając z regularności miary Lebesgue'a na zbiorze
otwartym $\{M\nu>t\}$, przechodzimy do supremum po
jego zwartych podzbiorach. W przypadku miary zespolonej
lub ze znakiem stosujemy otrzymany wynik do $|\nu|$.
\end{proof}

\begin{uwaga}[Nie mylić dwóch operatorów maksymalnych]
Dla $f\in L^1(\mathbb R^d)$ miara
$\nu_f(A)=\int_A|f(y)|\,\mathrm d\lambda_d(y)$ spełnia
$M\nu_f=Mf$, gdzie $Mf$ zdefiniowano w
\ref{def:HL-max-24}. Operator $M\nu$ pozwala jednak
również badać miary osobliwe, nieposiadające gęstości.
\end{uwaga}

\section{Twierdzenie Lebesgue'a o różniczkowaniu całek}
\begin{lemat}[Przybliżanie lokalne funkcją ciągłą]
\label{lem:continuous-dense-25}
Dla dowolnego $f\in L^1(\mathbb R^d)$ i
$\varepsilon>0$ istnieje $\varphi\in C_c(\mathbb R^d)$
taka, że $\|f-\varphi\|_1<\varepsilon$.
\end{lemat}
\begin{proof}
Jest to szczególny przypadek twierdzenia o gęstości
$C_c^\infty(\mathbb R^d)$ w $L^1$, dowiedzionego w
rozdziale~\ref{ch:regularnosc-funkcji}. W szczególności nie
zakładamy tutaj różniczkowania całek przy konstrukcji
aproksymacji.
\end{proof}

\begin{twierdzenie}[Lebesgue'a o różniczkowaniu]
\label{thm:lebesgue-diff-25}
Jeżeli $f\in L^1_{\mathrm{loc}}(\mathbb R^d)$, to
prawie każdy $x\in\mathbb R^d$ jest punktem Lebesgue'a
$f$, tzn. spełnia \eqref{eq:lebesgue-point-25}.
W szczególności
$\lim_{r\downarrow0}A_rf(x)=f(x)$ dla prawie każdego $x$.
\end{twierdzenie}
\begin{proof}
Najpierw załóżmy $f\in L^1(\mathbb R^d)$. Ustalmy
$\varphi\in C_c(\mathbb R^d)$ i połóżmy $g=f-\varphi$.
Dla każdego $x$, dla którego $f(x)$ jest określona,
z nierówności trójkąta otrzymujemy
\begin{align*}
 \limsup_{r\downarrow0}\frac1{v_dr^d}
 \int_{B(x,r)}|f(y)-f(x)|\,\mathrm d\lambda_d(y)
 &\le Mg(x)+|g(x)|,
\end{align*}
gdzie $Mg=M(|g|\lambda_d)$. Składnik z funkcją
$\varphi$ znika wskutek jej ciągłości w $x$.
Niech $D_f(x)$ oznacza powyższy $\limsup$, a $t>0$.
Dla każdego $\varphi$ zachodzi inkluzja
$\{D_f>2t\}\subseteq\{Mg>t\}\cup\{|g|>t\}$.
Nierówność słabego typu i nierówność Czebyszewa dają
\[
 \lambda_d(\{D_f>2t\})
 \le\frac{5^d+1}{t}\|f-\varphi\|_1.
\]
Lemat~\ref{lem:continuous-dense-25} pozwala uczynić
prawą stronę dowolnie małą. Zatem
$\lambda_d(\{D_f>2t\})=0$ dla każdego $t>0$.
Ponieważ $D_f\ge0$, otrzymujemy $D_f=0$ prawie
wszędzie. (Mierzalność $D_f$ wynika z mierzalności
średnich i z wyrażenia granicy górnej przy promieniach
wymiernych, wykorzystując ciągłość całki względem
promienia.)

Jeżeli $f$ jest tylko lokalnie całkowalna, dla każdego
$n\ge1$ rozważamy $f_n=f\chi_{B(0,n+2)}\in L^1$.
Dla $x\in B(0,n)$ i $r<1$ średnie funkcji
$f_n$ i $f$ są jednakowe, a wartości obu funkcji
w $x$ są równe. Zastosowanie pierwszej części do
wszystkich $n$ i przeliczalna suma zbiorów zerowych
dają tezę.
\end{proof}

\begin{wniosek}[Wartość funkcji z otaczających ją średnich]
\label{cor:diff-basis-25}
Jeżeli $f\in L^1_{\mathrm{loc}}(\mathbb R^d)$,
to istnieje zbiór $N$ o $\lambda_d(N)=0$, taki że
$\lim_{r\downarrow0}A_rf(x)=f(x)$ dla wszystkich
$x\in\mathbb R^d\setminus N$. Zbioru wyjątkowego
nie można w ogólności usunąć.
\end{wniosek}
\begin{proof}
Pierwsza część wynika z twierdzenia
\ref{thm:lebesgue-diff-25}; druga z przykładu
$f=\chi_{\mathbb Q^d}$.
\end{proof}

\begin{przyklad}[Przedstawiciel kanoniczny w punktach Lebesgue'a]
Jeżeli $f\in L^1_{\mathrm{loc}}(\mathbb R^d)$, to
na zbiorze punktów, w których istnieje skończona granica
$\lim_{r\downarrow0}A_rf(x)$, można ją przyjąć za
wartość przedstawiciela $\widetilde f(x)$. W pozostałych
punktach przyjmujemy $0$. Otrzymujemy funkcję mierzalną
$\widetilde f=f$ prawie wszędzie. Zmiana przedstawiciela
$f$ na zbiorze zerowym nie zmienia wartości żadnej średniej,
więc $\widetilde f$ nie zależy od tego wyboru.
\end{przyklad}

\section{Różniczkowanie miar: część absolutnie ciągła i osobliwa}
\begin{lemat}[Zanik gęstości miary osobliwej]
\label{lem:singular-density-25}
Niech $\nu$ będzie skończoną dodatnią miarą borelowską
na $\mathbb R^d$ i $\nu\perp\lambda_d$. Wówczas
\begin{equation}\label{eq:singular-density-25}
 \lim_{r\downarrow0}\frac{\nu(B(x,r))}{v_dr^d}=0
 \quad\text{dla }\lambda_d\text{-prawie każdego }x.
\end{equation}
\end{lemat}
\begin{proof}
Wybierzmy borelowski $N$ taki, że $\lambda_d(N)=0$
i $\nu(\mathbb R^d\setminus N)=0$. Regularność
skończonej miary borelowskiej daje dla każdego
$\varepsilon>0$ zwarty $K\subseteq N$ taki, że
$\nu(\mathbb R^d\setminus K)<\varepsilon$.
Oznaczmy $\nu_0=\nu|_K$, $\nu_1=\nu|_{\mathbb R^d\setminus K}$.
Dla każdego $x\notin K$ odległość $\operatorname{dist}(x,K)$
jest dodatnia, zatem $\nu_0(B(x,r))=0$ dla dostatecznie
małych $r$. W konsekwencji, poza $K$,
\[
 \limsup_{r\downarrow0}\frac{\nu(B(x,r))}{v_dr^d}
 \le M\nu_1(x).
\]
Niech $t>0$. Z twierdzenia~\ref{thm:weak-measure-25}
i z $\lambda_d(K)=0$ otrzymujemy
\[
 \lambda_d\left(\left\{x:\limsup_{r\downarrow0}
 \frac{\nu(B(x,r))}{v_dr^d}>t\right\}\right)
 \le\frac{5^d\varepsilon}{t}.
\]
Ponieważ $\varepsilon$ jest dowolne, zbiór ten ma
miarę zero. Dla $t=1/n$, $n\ge1$, dostajemy
zerowanie granicy górnej prawie wszędzie. Ilorazy są
nieujemne, więc istnieje granica równa zero.
\end{proof}

\begin{twierdzenie}[Różniczkowanie miary względem Lebesgue'a]
\label{thm:measure-differentiation-25}
Niech $\nu$ będzie dodatnią miarą Radona na
$\mathbb R^d$. Jej rozkład Lebesgue'a względem
$\lambda_d$ ma postać $\nu=f\lambda_d+\nu_s$,
gdzie $f\ge0$ jest lokalnie całkowalna, a
$\nu_s\perp\lambda_d$. Wtedy
\begin{equation}\label{eq:measure-density-25}
 \lim_{r\downarrow0}\frac{\nu(B(x,r))}{v_dr^d}
 =f(x)\qquad\text{dla }\lambda_d\text{-prawie każdego }x.
\end{equation}
Dla lokalnie skończonych miar ze znakiem oraz miar
zespolonych twierdzenie pozostaje prawdziwe, jeżeli
$f$ jest odpowiednią pochodną Radona--Nikodyma.
\end{twierdzenie}
\begin{proof}
Z rozkładu Lebesgue'a i twierdzenia Radona--Nikodyma
(rozdziały~\ref{ch:ac-singular-17} i
\ref{ch:radon-nikodym-18}) otrzymujemy
$\nu(B)=\int_Bf(y)\,\mathrm d\lambda_d(y)+\nu_s(B)$.
Pierwszy składnik po podzieleniu przez $v_dr^d$
zmierza do $f(x)$ prawie wszędzie z twierdzenia
\ref{thm:lebesgue-diff-25}. Dla składnika osobliwego
lokalizujemy miarę na kulach $B(0,n)$ i stosujemy
lemat~\ref{lem:singular-density-25}; przy
$x\in B(0,n-1)$ obie miary mają takie same wartości
na dostatecznie małych kulach. Przeliczalna suma
zbiorów wyjątkowych jest zerowa. W przypadku miar ze
znakiem i zespolonych stosujemy ten sam argument do
części absolutnie ciągłej, a do wariacji części osobliwej
stosujemy lemat~\ref{lem:singular-density-25}.
\end{proof}

\begin{uwaga}[Ważne rozróżnienie]
Istnienie granicy w \eqref{eq:measure-density-25} dla
$\lambda_d$-prawie wszystkich $x$ \emph{nie} oznacza,
że $\nu_s=0$. Dla miary Cantora pochodna gęstości
względem $\lambda_1$ jest zerowa prawie wszędzie,
chociaż sama miara jest niezerowa.
\end{uwaga}

\begin{przyklad}[Miara Diraca]
Niech $\nu=\delta_0$ na $\mathbb R^d$. Jeżeli $x\ne0$,
to dla wszystkich dostatecznie małych $r$ zachodzi
$\delta_0(B(x,r))=0$. Natomiast
$\delta_0(B(0,r))/(v_dr^d)=1/(v_dr^d)\to+\infty$.
Miara osobliwa może mieć więc nieskończoną gęstość
w swoich punktach koncentracji.
\end{przyklad}

\section{Pochodna całki w jednym wymiarze}
\begin{twierdzenie}[Podstawowe twierdzenie rachunku całkowego dla całki Lebesgue'a]
\label{thm:FTC-lebesgue-25}
Niech $a<b$ i $f\in L^1((a,b),\lambda_1)$. Funkcja
$F:[a,b]\to\mathbb R$ określona wzorem
$F(x)=\int_{(a,x)}f(t)\,\mathrm d\lambda_1(t)$
jest absolutnie ciągła i spełnia
$F'(x)=f(x)$ dla prawie każdego $x\in(a,b)$.
\end{twierdzenie}
\begin{proof}
Dla skończonej rodziny parami rozłącznych przedziałów
$(a_j,b_j)\subseteq[a,b]$ mamy
$\sum_j|F(b_j)-F(a_j)|
\le\int_{\cup_j(a_j,b_j)}|f|\,\mathrm d\lambda_1$.
Bezwzględna ciągłość całki zapewnia, że dla danego
$\varepsilon>0$ prawa strona jest mniejsza od
$\varepsilon$, gdy suma długości przedziałów jest
dostatecznie mała. Jest to definicja absolutnej
ciągłości funkcji $F$.

Niech $x\in(a,b)$ będzie punktem Lebesgue'a $f$.
Dla $h>0$ tak małego, by $x+h<b$, zachodzi
\[
 \left|\frac{F(x+h)-F(x)}h-f(x)\right|
 \le\frac1h\int_{(x,x+h)}|f(t)-f(x)|\,
 \mathrm d\lambda_1(t)
 \le\frac2{2h}\int_{(x-h,x+h)}|f(t)-f(x)|\,
 \mathrm d\lambda_1(t).
\]
Ostatnia średnia zmierza do zera. Analogiczny argument
obowiązuje dla $h<0$. Prawie wszystkie punkty są
punktami Lebesgue'a na mocy twierdzenia
\ref{thm:lebesgue-diff-25}, zatem $F'=f$ prawie
wszędzie.
\end{proof}

\begin{przyklad}[Dlaczego sama pochodna prawie wszędzie nie wystarcza?]
Klasyczna funkcja Cantora $C:[0,1]\to[0,1]$ jest
ciągła i niemalejąca, a $C'=0$ prawie wszędzie, lecz
$C(1)-C(0)=1$. Nie jest absolutnie ciągła.
Dlatego z samej równości $F'=f$ prawie wszędzie nie
wynika $F(b)-F(a)=\int_{(a,b)}f\,\mathrm d\lambda_1$.
\end{przyklad}

\begin{twierdzenie}[Funkcja rozkładu a absolutna ciągłość miary]
\label{thm:cdf-derivative-25}
Niech $\nu$ będzie skończoną dodatnią miarą borelowską
na $\mathbb R$. Określmy $F:\mathbb R\to[0,\nu(\mathbb R)]$
wzorem $F(x)=\nu(( -\infty,x])$.
Jeśli $\nu=f\lambda_1+\nu_s$ jest rozkładem Lebesgue'a
miary $\nu$ względem $\lambda_1$, to
$F'(x)=f(x)$ dla $\lambda_1$-prawie każdego $x$.
Ponadto następujące warunki są równoważne:
\begin{enumerate}[label=(\roman*)]
\item $\nu\ll\lambda_1$;
\item $F$ jest absolutnie ciągła na każdym zwartym
przedziale;
\item dla wszystkich $a<b$ zachodzi
$F(b)-F(a)=\int_{(a,b)}F'(t)\,\mathrm d\lambda_1(t)$.
\end{enumerate}
W (iii) pochodna istnieje prawie wszędzie.
\end{twierdzenie}
\begin{proof}
Niech $x$ będzie punktem Lebesgue'a $f$, w którym
$\nu_s(B(x,r))/(2r)\to0$ oraz $\nu(\{x\})=0$.
Punkty takie tworzą zbiór pełnej miary Lebesgue'a:
pierwsze dwie własności wynikają z twierdzenia
\ref{thm:measure-differentiation-25}, a atomów dodatniej
masy jest co najwyżej przeliczalnie wiele.
Dla $h>0$ mamy
\[
 \frac{F(x+h)-F(x)}h
 =\frac1h\int_{(x,x+h]}f(t)\,\mathrm d\lambda_1(t)
 +\frac{\nu_s((x,x+h])}h.
\]
Pierwszy składnik dąży do $f(x)$, ponieważ średnia
wartości $|f(t)-f(x)|$ po $(x,x+h)$ jest ograniczona
przez dwa razy analogiczną średnią po $(x-h,x+h)$.
Drugi składnik, dla $h>0$, jest ograniczony przez
$\nu_s(B(x,2h))/h$ i zatem dąży do zera.
Dla $h<0$ rozpatrujemy przedział $(x+h,x]$;
ponieważ $\nu_s(\{x\})=0$, ten sam argument działa
również z lewej. Dowodzi to $F'=f$ prawie wszędzie.

Jeśli (i) zachodzi, to $\nu=f\lambda_1$ z
$f\in L^1(\mathbb R)$, a więc
$F(x)=\int_{(-\infty,x]}f(t)\,\mathrm d\lambda_1(t)$.
Bezwzględna ciągłość całki na zbiorach o małej mierze
daje (ii), zaś podstawowe twierdzenie rachunku
całkowego (twierdzenie~\ref{thm:FTC-lebesgue-25})
daje (iii).

Załóżmy (ii) i wykażmy najpierw (i), bez powoływania
się na twierdzenie o całkowaniu pochodnej dowolnej funkcji
absolutnie ciągłej. Ustalmy $n\ge1$ i zbiór borelowski
$E\subseteq[-n,n]$ z $\lambda_1(E)=0$. Funkcja $F$
jest absolutnie ciągła na $[-n-1,n+1]$, więc dla
każdego $\varepsilon>0$ istnieje $\delta>0$ takie,
że dla każdej skończonej rodziny rozłącznych przedziałów
zawartych w $[-n-1,n+1]$, o łącznej długości poniżej
$\delta$, suma przyrostów $F$ jest mniejsza od
$\varepsilon$. Pokryjmy $E$ otwartym zbiorem
$O\subset(-n-1,n+1)$ o $\lambda_1(O)<\delta$.
Składowymi $O$ są parami rozłączne przedziały otwarte
$(a_j,b_j)$. Ponieważ $F$ jest ciągła,
$\nu(\{x\})=F(x)-F(x^-)=0$ dla każdego $x$
wewnątrz $[-n-1,n+1]$, a więc
$\nu((a_j,b_j))=F(b_j)-F(a_j)$.
Stosując nierówność absolutnej ciągłości do skończonych
sum i przechodząc do granicy z przeliczalnej addytywności,
dostajemy $\nu(E)\le\nu(O)\le\varepsilon$.
Zatem $\nu(E)=0$. Ponieważ $n$ był dowolny,
$\nu\ll\lambda_1$ na całej prostej. Wykazane już
(i)$\Rightarrow$(iii) daje również (ii)$\Rightarrow$(iii).
Załóżmy wreszcie (iii). Ponieważ
$\nu((a,b])=F(b)-F(a)$, dla każdych $a<b$ mamy
$\nu((a,b])=\int_{(a,b)}F'(t)\,\mathrm d\lambda_1(t)$;
przy całce nie ma znaczenia wybór końców.
Funkcja $F'\ge0$ prawie wszędzie i jest lokalnie
całkowalna, gdyż z (iii) jej całki po przedziałach
są skończone. Obie miary są skończone na zbiorach
ograniczonych i zgadzają się na przedziałach $(a,b]$.
Jednoznaczność miary na generującym $\pi$-układzie
i przejście przez zbiory $(-n,n]$ pokazują
$\nu(E)=\int_EF'(t)\,\mathrm d\lambda_1(t)$
dla wszystkich $E\in\mathcal B(\mathbb R)$.
To daje (i).
\end{proof}

\begin{przyklad}[Pochodna nie wykrywa składnika osobliwego]
W przypadku miary Cantora funkcja rozkładu $F$ jest
ciągła, niemalejąca i spełnia $F'=0$ prawie wszędzie,
ale nie jest absolutnie ciągła. Twierdzenie
\ref{thm:cdf-derivative-25} wyjaśnia przyczynę:
pochodna odzyskuje wyłącznie część absolutnie ciągłą
miary, a nie składnik osobliwy.
\end{przyklad}

\section{Przenoszenie różniczkowania na inne rodziny zbiorów}
\begin{wniosek}[Średnie po kostkach]
\label{cor:cubes-differentiation-25}
Niech $f\in L^1_{\mathrm{loc}}(\mathbb R^d)$. Dla prawie
każdego $x$ średnia $f$ po dowolnych kostkach osiowych
$Q_r(x)=x+(-r,r)^d$ dąży do $f(x)$ przy $r\downarrow0$.
Ponadto $\lambda_d(Q_r(x))^{-1}
\int_{Q_r(x)}|f(y)-f(x)|\,\mathrm d\lambda_d(y)\to0$.
\end{wniosek}
\begin{proof}
Mamy $Q_r(x)\subseteq B(x,\sqrt d\,r)$ oraz
$\lambda_d(B(x,\sqrt d\,r))/\lambda_d(Q_r(x))
=v_dd^{d/2}/2^d$, niezależnie od $r$. Dodatnia średnia
po kostce jest więc ograniczona przez stałą razy
średnia po większej kuli, która zmierza do zera
w każdym punkcie Lebesgue'a.
\end{proof}

\begin{uwaga}[Przekroje i inne miary]
Analogiczne różniczkowanie względem \emph{dowolnej}
miary Radona $\mu$ na $\mathbb R^d$ można otrzymać
z twierdzenia pokryciowego Besicovitcha:
centrowany operator maksymalny względem $\mu$ spełnia
odpowiednie oszacowanie słabego typu $(1,1)$, a dalszy
argument opiera się na aproksymacji funkcji całkowalnych.
Pełna, nieskończona wersja twierdzenia Besicovitcha
\emph{została już udowodniona} w
Twierdzeniu~\ref{thm:besicovitch-24}.
Nie wyprowadzamy tutaj jednak osobno całego twierdzenia
o różniczkowaniu względem arbitralnej miary Radona:
wymaga ono dodatkowo sformułowania i przeprowadzenia
argumentu maksymalnego z mianownikiem $\mu(B(x,r))$,
z uwzględnieniem kul o zerowej mierze.
W ten sposób rozróżniamy udowodnione narzędzie pokryciowe
od niewypisanego w pełni wniosku o różniczkowaniu.
\end{uwaga}